\chapter{认识保研：制度、路径与时间线}
\label{ch01}
\ist{认知与规划}
\begin{keybox}
本章讲清三件事：推免到底是什么、三年时间怎么排、信息从哪来。

最关键的三条全国性规则：第一，推免免掉的只是全国统考初试，复试、报名手续、公示备案一样不少（教学〔2006〕14号）。第二，名额由教育部按学校层次核定并下达——设研究生院的高校一般按应届本科毕业生数的15\%左右、未设研究生院的“211工程”高校约5\%、其他高校约2\%，学校再切到学院和专业。第三，2014年起推免名额不再区分学术学位与专业学位、不再设留校限额，推免生享有自主选择报考招生单位和专业的权利。

最易踩的四个坑：把填志愿当第一步（其实要先在校内拿到推免资格，北京大学2027届9月1日至7日就开放资格申请）；漏看省级教育考试机构公告（2027届9月18日注册、9月21日9:00开始填报志愿、10月18日截止填报、10月20日前未落实接收单位不再保留资格）；以为夏令营拿到意向就稳了（未经过教育部推免服务系统接收的推免生一律无效）；不知道自己被排除（专升本、第二学士学位、独立学院学生全国统一不在推荐对象范围内）。
\end{keybox}

\section{公平性}
\begin{qa}{M0002}{保研和考研到底差在哪？听说保研能省一整年是真的吗}
\textbf{不是。推免免掉的只是全国统考的初试，不是整个读研过程；省一年是夸大。}

教学〔2006〕14号第三条：推免生不必经过全国硕士研究生入学统一考试的初试，直接进入招生单位复试。差别有三：

\begin{itemize}
  \item \textbf{考核方式}：推免看本科前三年学业成绩与校内综合评价，再去参加目标院校复试；统考拼12月初试分数。
  \item \textbf{结果时点}：2027届推免志愿填报在2026年9月21日至10月18日，待录取确认9月24日开始（四川省教育考试院2027年公告）；统考要等12月考试、次年春季录取。
  \item \textbf{入学一致}：推免生同样次年9月入学，大四照样要修完课程、取得毕业证与学位证。
\end{itemize}

两条路不可兼得：已录取的推免生不得再报名当年统考；在截止日期未被录取的推免生不再保留资格，可报考当年统考。
\nref{S0454}
\nref{S0389}
\nref{S0352}
\end{qa}

\iin{教育部学生服务与素质发展中心（学信网/研招网）}
\ith{考研对比}
\ist{认知与规划}
\begin{qa}{M0004}{大一高数挂过一次，是不是直接跟保研说再见了}
\textbf{全国文件没有挂科一票否决的条款，但多数学校细则把成绩优良写成硬条件。挂过科不等于出局，关键看本校当年细则。}

教学〔2006〕14号第十四条只提勤奋学习、刻苦钻研、成绩优秀这类原则要求，未规定某门课不及格即取消资格。决定权在校规：

\begin{itemize}
  \item \textbf{原则上不予推荐}：北京大学2027届办法写原则上必修课有不及格或学业成绩处于下游的不予推荐。
  \item \textbf{无不及格记录}：广东外语外贸大学2026年实施办法把无课程不及格记录列为学业表现基本要求之一。
  \item \textbf{给门数上限}：大连理工大学商学院规定一考不及格的必修课程不能超过2门。
  \item \textbf{影响录取环节}：浙江大学2027年通知明确，拟录取后本科最后一学年计入毕业要求学分的课程出现不及格记录，取消录取资格。
\end{itemize}

先查本校细则是否写无不及格记录或允许几门，再看能否把排名拉回来。
\nref{S0454}
\nref{S0553}
\nref{S0134}
\nref{S0549}
\nref{S0032}
\end{qa}

\iin{教育部学生服务与素质发展中心（学信网/研招网）}
\ith{挂科影响}
\ist{认知与规划}
\begin{qa}{M0010}{保研率一般多少算正常？我们专业2\%是不是太低了}
\textbf{官方文件里没有保研率这个词，只有推免名额。比例由国家按学校层次定原则、学校再往下切，2\%本身判断不了高低。}

教学〔2006〕14号第十二条的名额原则：

\par\vspace{2pt}\nobreak\begin{center}\footnotesize
\begin{tabularx}{\textwidth}{>{\raggedright\arraybackslash}X>{\raggedright\arraybackslash}X}
\toprule
\textbf{学校类型} & \textbf{名额基准} \\ \midrule
设研究生院的高校 & 应届本科毕业生数的15\%左右 \\
未设研究生院的211工程高校 & 5\%左右 \\
其他高校 & 2\%，初次开展工作的前3年按1\%左右 \\
人文、理科等人才培养基地班 & 按基地班招生人数的50\%左右单独增加 \\
\bottomrule
\end{tabularx}
\end{center}\vspace{2pt}\par

这只是学校层面的基准。名额到校后还要二次分配：湖南大学2027年通知写明学校按教育部下达名额分配至有应届毕业生的学院；广东外语外贸大学2026年办法综合考虑毕业生人数、硕士点、专业数量与上年度名额使用情况。切到专业后2\%甚至更低都可能出现。
\nref{S0454}
\nref{S0540}
\nref{S0134}
\end{qa}

\iin{教育部学生服务与素质发展中心（学信网/研招网）}
\ith{保研率}
\ist{认知与规划}
\begin{qa}{M0012}{直博是啥？保研选了直博是不是就退不出来了}
\textbf{直博是从推免生中选拔、直接取得博士入学资格的方式，基本学制通常5年。退不出来不准确，但退出成本高、有明确时间限制。}

中山大学2027年推免办法写明直博生学制为5年；中南大学2027年选拔推免生直接攻博的通知同样写学制5年，并给出关键出口：

\begin{itemize}
  \item \textbf{可转硕士层次}：按《中南大学研究生培养环节工作管理办法》第八条，第八学期前未通过资格考试的直博生可转入硕士层次继续学习或作退学处理。
  \item \textbf{放弃有截止日}：拟录取直博生如因故放弃，须在2027年3月15日前向研招办出具书面报告；\textbf{放弃资格后不能转录为硕士研究生}。
  \item \textbf{取消资格情形}：材料与事实不符、本科最后一年有不及格科目、受纪律处分、未取得毕业证与学位证等。
\end{itemize}

报考前务必先读目标院校的研究生培养管理规定，确认资格考试时间点和分流条件。
\nref{S0167}
\nref{S0625}
\nref{S0635}
\end{qa}

\iin{中山大学}
\ith{直博认知}
\ist{认知与规划}
\begin{qa}{M0018}{家里人说保研比考研稳，是这样吗}
\textbf{要分开看。制度确定性上推免比统考强，但前提是先拿到推免资格，这个前置门槛本身就不确定。}

推免的确定性来源：

\begin{enumerate}
  \item \textbf{时间早}：2027届志愿填报在2026年9月21日至10月18日，10月20日前未落实接收单位不再保留资格（四川省教育考试院2027年公告），比统考早出结果约半年。
  \item \textbf{免初试}：不必参加全国统考初试，直接进复试（教学〔2006〕14号第三条）。
  \item \textbf{程序受约束}：名额由教育部下达，学校按综合评价择优推荐，名单公示期不少于7天，未经公示的推免生资格无效。
\end{enumerate}

不确定性集中在两处：一是\textbf{资格关}，名额切到学院专业后可能只有几个人；二是\textbf{接收关}，拿到资格不等于有学校要。统考的不确定性则集中在一次考试，两种不稳性质不同。
\nref{S0454}
\nref{S0389}
\nref{S0352}
\end{qa}

\iin{教育部学生服务与素质发展中心（学信网/研招网）}
\ith{考研对比}
\ist{认知与规划}
\begin{qa}{M0026}{少数民族骨干计划算保研名额吗}
\textbf{不算。少数民族高层次骨干人才计划是与推免并列的另一条招生通道，硕士层面主要走统考。}

贵州省教育厅2026年骨干计划报考通知（转引教育部教民厅〔2025〕1号）明确，骨干计划硕士实行自愿报考、统一考试、单独划线、择优录取，由教育部按学科门类分别确定进入复试的初试成绩基本要求。主体通道是统考。

两者的交集是：推免生\textbf{可以}申请骨干计划。中国人民大学2025年骨干计划招生简章在硕士报考条件里写明符合条件的推荐免试生可申请本计划。山东省教育招生考试院2026年说明指出，报考者须通过研招网骨干计划资格申请平台提交申请、审核通过后方可报名。

两条容易踩的坑：报考骨干计划硕士不再享受其他初试加分和照顾政策，且\textbf{不得兼报其他专项计划或硕士研究生普通计划}；录取后须签定向就业协议，毕业回定向地区单位就业。
\nref{S0573}
\nref{S0578}
\nref{S0595}
\end{qa}

\iin{中国政府网}
\ith{少数民族}
\ist{认知与规划}
\begin{qa}{M0040}{保研是不是更看关系背景}
\textbf{从制度设计看，推免套着比多数考试招生更密的程序约束；但完全看关系和完全无关都不成立，中间地带取决于具体学校的执行。}

可查到的约束：

\begin{itemize}
  \item \textbf{标准与程序须公开}：推免工作应做到公开、公正、公平，推荐学校与接收单位均须制订科学、规范、明确的推荐标准、接收标准及公开透明的工作程序（教学〔2006〕14号第五条）。
  \item \textbf{结果须公示}：名单在各院（系）和校内网站公示，公示期不少于7天；有异议须查明情况、公布处理结果；未经公示的推免生资格无效（第十五条）。
  \item \textbf{申诉进校内渠道}：高校应将推免工作中的学生申诉纳入校内申诉渠道（第二十三条）。
  \item \textbf{回避制度}：北京大学2027届办法要求中层及以上领导干部和推免工作人员如有直系亲属或其他利益相关者报名，须主动报备并全程回避。
\end{itemize}

这些条款不等于执行到位，但给了你投诉和申辩的抓手。
\nref{S0454}
\nref{S0553}
\nref{S0568}
\end{qa}

\iin{教育部学生服务与素质发展中心（学信网/研招网）}
\ith{公平性}
\ist{认知与规划}
\begin{qa}{M0041}{保研名额会不会被内定？公示只是走个过场吗}
\textbf{公示不是走过场。全国文件明确规定未经公示的推免生资格无效；对公示有异议，学校必须查明并公布处理结果。}

教学〔2006〕14号第十五条：学校将审定的名单在各院（系）和校内网站公示，\textbf{公示期不少于7天}；对有异议的学生，学校要查明情况、公布处理结果；无异议的报省级招办备案；未经公示的推免生资格无效。第二十二条还要求公开推免政策规定、推免生资格、招收名额、录取信息、咨询申诉渠道。

不同学校分层设置公示，例如：

\begin{itemize}
  \item 中南大学材料科学与工程学院2027届预推荐名单公示时间为2026年9月13日至9月19日共七天，并公布书面反映渠道、电话和邮箱。
  \item 青海大学的流程是院（系）公示3天，报教务处后全校公示10天，期满无异议或异议不成立，经校党委会审议通过后确定推荐名单。
\end{itemize}

有疑问时，拿本院当年细则与公示名单核对综合成绩计算是否一致，再按公示渠道书面提异议。
\nref{S0454}
\nref{S0274}
\nref{S0232}
\end{qa}

\iin{教育部学生服务与素质发展中心（学信网/研招网）}
\ith{名额公示}
\ist{认知与规划}
\begin{qa}{M0042}{我想跨专业保研，是不是太天真了}
\textbf{跨专业保研不是天真，全国文件本身鼓励学科交叉；难点在接收院系是否愿意收、以及你要补多少专业课。}

制度依据有两层：

\begin{itemize}
  \item \textbf{推荐端鼓励流动}：教学〔2006〕14号第七条提倡优势互补、加强交流，鼓励推免生在不同地区间、招生单位间及不同学科间的交流；北京大学2027届办法也要求各院系积极推荐本校毕业生到其他院校、科研机构攻读研究生，不囿于本校、本院系、本专业。
  \item \textbf{接收端由院系自定}：东北师范大学2027年推免办法明确，考生可申请跨学科攻读研究生，是否接收跨学科考生及具体考核要求由各招生院部自定。
\end{itemize}

操作上注意三点：原专业的排名绩点是推荐端基础，跨专业不降低拿资格的门槛；接收端考核通常含专业课，往往需提前自学目标专业核心课程；部分专业对本科背景有硬性要求，需逐个查目标院校推免专业目录。
\nref{S0454}
\nref{S0553}
\nref{S0607}
\end{qa}

\iin{教育部学生服务与素质发展中心（学信网/研招网）}
\ith{跨专业}
\ist{认知与规划}
\begin{qa}{M0047}{保研的人是不是都不用考研了}
\textbf{不对。推免免掉的只是全国统考初试，复试照样要参加；而且推免生被录取后仍要走一遍统考网上报名手续。}

教学〔2006〕14号第三条把免试定义得很清楚：指应届本科毕业生不必经过全国硕士研究生入学统一考试的初试，\textbf{直接进入复试}。同办法第十九条规定，具有推免生资格的考生被录取后仍须参加全国硕士研究生招生考试网上报名，并按规定到指定报考点现场确认；第二十条明确未经复试的申请者不得录取。

反方向也常被混淆：

\begin{itemize}
  \item 已录取的推免生不得再报名参加当年全国硕士研究生招生考试。
  \item 在规定截止日期未被招生单位录取的推免生不再保留推免资格，可在网上报名截止时间内报考当年统考。
\end{itemize}

所以那句话只在不用考初试这一个意义上成立。推免更接近提前批录取：少了初试，复试、报名手续、公示备案一个都不少。
\nref{S0454}
\nref{S0352}
\nref{S0389}
\end{qa}

\iin{教育部学生服务与素质发展中心（学信网/研招网）}
\ith{路径关系}
\ist{认知与规划}
\begin{qa}{M0055}{保研和高考比哪个更公平}
\textbf{两者无法用同一把尺子量。高考是全国统一考试、按分投档；推免是校内排序加校外考核的双重筛选，比较的其实是两种公平观。}

高考的公平性建立在统一命题、统一评分、按分录取的单一尺度上，争议点主要在省份间招生计划分配。推免的公平性建立在程序公开上：

\begin{itemize}
  \item 推荐端由学校按综合评价择优遴选，学业综合成绩是最基础的指标，可纳入参军入伍、志愿服务、国际组织实习、科研创新成果、竞赛奖项等，但\textbf{不得单列名额、单独排序}（北京大学2027届办法）。
  \item 名单公示期不少于7天，未经公示的推免生资格无效；申诉纳入校内申诉渠道（教学〔2006〕14号第十五条、第二十三条）。
  \item 接收端仍要向招生单位投递、复试，能否录取取决于对方院系判断。
\end{itemize}

可以确定的差异是：高考的不公平通常表现为区域名额差异，推免的不公平通常表现为院校出身和本校细则口径差异。
\nref{S0454}
\nref{S0553}
\nref{S0048}
\end{qa}

\iin{教育部学生服务与素质发展中心（学信网/研招网）}
\ith{公平性}
\ist{认知与规划}
\begin{qa}{M0057}{保研名额没用完会顺延给下一名吗}
\textbf{要分两种情形。名额被放弃后转让给别人是明确禁止的；学院内部的候补、递补是学校常见做法，两者不是一回事。}

\textbf{禁止的部分（全国统一）}：教学〔2006〕14号第十三条规定，高校不得将推免生名额跨学校使用；截至教育部规定时间后，被确定的推免生如果无招生单位接收或本人放弃，则该生推免生资格作废，\textbf{名额不得转让}。

\textbf{学校内部候补}：很多学校在遴选阶段就设候补名单补前序空缺。可查实例：

\begin{itemize}
  \item 成都大学商学院2027届细则规定候补人选人数按各专业推免生名额数的50\%计算（不足1人的按1人计）。
  \item 浙江师范大学2027年发布过《2027年免试攻读硕（博）士学位研究生递补名单公示》；中南大学材料科学与工程学院也公示过推免增补名单。
\end{itemize}

所以在候补位次上确实可能因前面有人放弃而补进来，但已生效的资格放弃后，空出名额不会被私下转给别人。
\nref{S0454}
\nref{S0079}
\nref{S0570}
\nref{S0275}
\end{qa}

\iin{教育部学生服务与素质发展中心（学信网/研招网）}
\ith{名额分配}
\ist{认知与规划}
\section{准备时间线}
\begin{qa}{M0003}{我现在大二下，绩点3.4还能保研吗？来得及补吗}
\textbf{全国文件没有绩点门槛，3.4够不够取决于本校换算方式和你的专业排名。大二下开始补，时间上够。}

全国层面只有原则性表述：教学〔2006〕14号第十四条要求成绩优秀，没给分数线。门槛由学校定：

\begin{itemize}
  \item 广东外语外贸大学2026年办法要求平均学分绩点不低于3.0，无课程不及格记录，并完成应修总学分的80\%。
  \item 北京大学2027届办法写原则上必修课有不及格或学业成绩处于下游的不予推荐。
  \item 大连理工大学商学院按学习成绩85\%、综合素质和能力测评10\%加权计算推免成绩，一考不及格的必修课程不能超过2门。
\end{itemize}

判断顺序：先查本校是否设绩点下限，再查本院名额数（湖南大学2027年通知写明学校按教育部下达名额分配到学院，学院再切到专业），然后看\textbf{专业排名}而非绝对绩点。大二下到大三上是绩点权重最大的阶段。
\nref{S0454}
\nref{S0134}
\nref{S0553}
\nref{S0549}
\nref{S0540}
\end{qa}

\iin{教育部学生服务与素质发展中心（学信网/研招网）}
\ith{准备时间线}
\ist{认知与规划}
\begin{qa}{M0013}{我准备考公，还能保研吗？会不会冲突}
\textbf{制度上不互斥，保研和考公可以同时报名。冲突在两处：时间重叠，以及推免录取后你就不能再按普通应届生的节奏拖到次年定去向。}

全国层面的硬约束：

\begin{itemize}
  \item 2027届推免志愿填报窗口是2026年9月21日至10月18日，待录取确认9月24日开始，10月20日前未落实接收单位则不再保留推免资格（四川省教育考试院2027年公告）。
  \item 已录取的推免生不得再报名参加当年全国硕士研究生招生考试，否则取消推免录取资格。
  \item 教学〔2006〕14号第二十条规定，入学前未取得学士学位或本科毕业证书，或受处分的，取消录取资格。
\end{itemize}

考公笔试时间不属本章教育口径，需查当年招考公告。可以确定的是：推免一旦确认待录取，你就进入次年9月入学的轨道，两条路都要投入，最后必须取舍。公费师范生、优师计划、农村订单定向生另有履约限制，中山大学2027年推免办法明确不接收未完成履约责任的这类推免生。
\nref{S0389}
\nref{S0352}
\nref{S0454}
\nref{S0167}
\end{qa}

\iin{www.sceea.cn}
\ith{路径冲突}
\ist{认知与规划}
\begin{qa}{M0015}{大三下才决定保研，还来得及吗？需要补什么}
\textbf{来得及，但要接受一个现实：夏令营第一批名额大概率错过，主战场要放在预推免和推免服务系统。}

三个时间锚点（2027届）：

\begin{enumerate}
  \item \textbf{校内资格申请}：北京大学2027届为2026年9月1日至7日开通校内推免资格申请系统，9月8日至10日院系初审，9月18日前学校审核公示。资格环节在大四开学初就收口。
  \item \textbf{夏令营}：集中在6—7月，报名截止常在6月中下旬，如北京大学力学与工程科学学院2026年夏令营要求6月22日08:30前完成材料上传。
  \item \textbf{推免服务系统}：2026年9月18日起注册填报个人资料并查看本人推免资格，9月21日9:00开始填报志愿，10月18日截止。
\end{enumerate}

大三下（3—6月）才决定，能抓住的是6月前后的夏令营末班车、8—9月预推免（如湖南大学2027年预报名8月28日起）、9月21日后的正式填报。要补的主要是成绩单与排名证明、英语成绩、能写进个人陈述的科研或实践经历。
\nref{S0553}
\nref{S0264}
\nref{S0389}
\nref{S0158}
\end{qa}

\iin{北京大学}
\ith{大三起步}
\ist{认知与规划}
\begin{qa}{M0019}{保研上岸的学长学姐一般什么时候开始准备的}
\textbf{没有统一答案，但从可查的官方节点倒推，最晚在大三下学期就必须进入实质准备。}

推免固定节点（2027届）：夏令营6—7月举办、报名截止多在6月中下旬；校内推免资格申请在大四开学初，北京大学2027届为9月1日至7日；推免服务系统9月18日起注册、9月21日9:00起填报志愿、10月18日截止（四川省教育考试院2027年公告）。倒推：

\begin{itemize}
  \item \textbf{大三下（3—6月）}：备材料、考英语、定目标院校，投夏令营。
  \item \textbf{大三上及以前}：维持学业排名，积累科研或竞赛经历。
\end{itemize}

需要注意：教学〔2006〕14号把推免定义为对优秀应届本科毕业生的遴选，北京大学2027届办法进一步写明把学生本科应届毕业学年前各学年的学业综合成绩作为遴选推荐综合评价体系的基础指标。决定你能不能被推荐的是前三年持续表现，排名从大一就在计分。
\nref{S0389}
\nref{S0553}
\nref{S0264}
\nref{S0454}
\end{qa}

\iin{www.sceea.cn}
\ith{准备时间线}
\ist{认知与规划}
\begin{qa}{M0023}{我是中外合作办学的，我们这种能保研吗}
\textbf{不能一概而论。中外合作办学学生能否推免，取决于该项目是否属于纳入国家普通本科招生计划录取，学校之间口径差异很大。}

全国层面依据是教学〔2006〕14号第十四条第（一）项：推荐对象为纳入国家普通本科招生计划录取的应届毕业生（不含专升本、第二学士学位、独立学院学生）。该条未直接点中外合作办学，判断权落在学校细则上。可以看到两种相反口径：

\textbf{可能在范围内}：浙江师范大学2027年免试攻读硕（博）士学位研究生拟推荐名单公示中，明确出现小学教育（中外合作办学）等专业的推荐学生。

\textbf{可能被排除}：广东外语外贸大学2026年办法限定为全日制应届本科毕业生，明确不含本硕项目的国际班学生；云南师范大学2027年通知的推荐对象也明确不含3+2高本贯通等学生。

此外，哈尔滨工业大学2027年接收工作相关通知把中外合作办学研究生、内地与港澳台地区合作办学研究生的学费标准单列。找本院当年细则看推荐对象那一栏最直接。
\nref{S0454}
\nref{S0571}
\nref{S0134}
\nref{S0568}
\nref{S0067}
\end{qa}

\iin{教育部学生服务与素质发展中心（学信网/研招网）}
\ith{中外合作}
\ist{认知与规划}
\begin{qa}{M0024}{专升本的同学有机会保研吗？听说名额特别少}
\textbf{不能。专升本学生被全国统一文件明确排除在推免推荐对象之外。}

教学〔2006〕14号第十四条第（一）项规定，推免生从纳入国家普通本科招生计划录取的应届毕业生（\textbf{不含专升本、第二学士学位、独立学院学生}）中择优遴选。这不是某一所学校的自选动作，而是全国层面写死的排除项。

此后多校在不同年份的细则里原文重申，可互相印证：

\begin{itemize}
  \item 北京大学2027届办法：申请条件第（1）项为纳入国家普通本科招生计划录取的应届毕业生（不含专升本、第二学士学位学生）。
  \item 广东外语外贸大学2026年办法：推荐对象不含辅修学位、专升本、第二学士学位学生、本硕项目的国际班学生和独立学院学生。
  \item 武汉轻工大学2027年接收推免生工作办法：申请条件明确不含专升本、第二学士学位学生。
\end{itemize}

被推免通道排除不等于不能读研：取得本科毕业证和学士学位后，可按普通应届本科毕业生身份报考全国硕士研究生招生考试。
\nref{S0454}
\nref{S0553}
\nref{S0134}
\nref{S0409}
\end{qa}

\iin{教育部学生服务与素质发展中心（学信网/研招网）}
\ith{专升本}
\ist{认知与规划}
\begin{qa}{M0027}{公费师范生能保研吗？听说要履行协议}
\textbf{公费师范生不走普通推免通道，走的是本研衔接专门通道，且必须履行定向任教协议。}

政策依据是《教育部直属师范大学本研衔接师范生公费教育实施办法》（国办发〔2024〕27号）。要点：从2024年起，国家在六所教育部直属师范大学实施本研衔接师范生公费教育（本科4年、教育硕士2年），支持符合条件的公费师范生\textbf{免试攻读本校全日制教育硕士研究生再履约任教}；研究生一年级课程学习结束后，按本科以来的综合考核结果排序，在履约任教地范围内选择任教地，毕业后从事中小学教育工作\textbf{6年以上}。

各校据此发布专门办法，如华中师范大学2027年公费师范生免试攻读教育硕士招生办法、东北师范大学2027年公费师范毕业生申请免试攻读非全日制教育硕士专业学位研究生招生办法。

另一面是限制：云南师范大学2027年通知的推荐对象明确不含公费师范生、优师计划师范生；中山大学2027年办法写明不接收未按国家规定完成履约责任（含违约）的农村订单定向生、教育部直属师范院校公费师范生等推免生申请。
\nref{S0456}
\nref{S0355}
\nref{S0559}
\nref{S0609}
\nref{S0568}
\nref{S0167}
\end{qa}

\iin{教育部学生服务与素质发展中心（学信网/研招网）}
\ith{公费师范}
\ist{认知与规划}
\begin{qa}{M0032}{我现在大一，想保研，这三年应该怎么规划}
\textbf{先把排名做扎实，再谈科研和竞赛。推免遴选的基础指标是学业综合成绩，前三年每年都在计分。}

教学〔2006〕14号第六条要求在对平时学习和科研能力综合测评的基础上，突出对创新精神、创新能力和专业能力倾向的考查。落到学校：北京大学2027届办法把学生本科应届毕业学年前各学年的学业综合成绩作为遴选推荐综合评价体系的基础指标，可纳入参军入伍服兵役、院系认可的志愿服务、到国际组织实习、高水平科研创新成果、竞赛奖项等，但\textbf{不得单列名额、单独排序}。

按三年倒推：

\begin{itemize}
  \item \textbf{大一}：打好高数、专业基础课和英语底子，避免不及格记录；了解本院推免细则的排名口径与名额数。
  \item \textbf{大二}：进入课题组或科创项目，把科研或竞赛经历从零变成一；确认英语成绩，六级425分以上是多校常见门槛。
  \item \textbf{大三}：3—6月投夏令营，6—7月参加，8—9月走预推免，9月21日推免服务系统开放后完成志愿填报。
\end{itemize}
\nref{S0454}
\nref{S0553}
\nref{S0134}
\end{qa}

\iin{教育部学生服务与素质发展中心（学信网/研招网）}
\ith{长期规划}
\ist{认知与规划}
\begin{qa}{M0035}{保研夏令营大概什么时候开始投？会不会撞期末}
\textbf{夏令营报名高峰是5—6月，举办高峰是6月底到7月中下旬，报名截止时间正好压在期末考试前后。}

从2026年各校通知可见时间分布：

\begin{itemize}
  \item 北京大学力学与工程科学学院2026年夏令营要求材料6月22日08:30前上传系统，纸质版6月22日16:30前签收。
  \item 山东大学经济研究院2026年夏令营报名截止7月1日中午12:00，活动安排在7月13日起。
  \item 浙江大学传媒与国际文化学院2026年夏令营于7月7—8日举行。
  \item 华中科技大学网络空间安全学院2026年夏令营于7月9—10日举办。
  \item 浙江大学生命科学研究院2026年夏令营定于7月14—17日举办。
\end{itemize}

报名截止集中在6月中下旬至7月初，恰好是多数学校的期末考试周。应对办法是提前备材料：成绩单、排名证明、英语成绩、推荐信在5月之前就该备好。

错过夏令营不等于结束：2027届推免服务系统填报志愿要到2026年9月21日才开始，10月18日截止，预推免和正式填报仍有大量机会。
\nref{S0264}
\nref{S0013}
\nref{S0003}
\nref{S0389}
\nref{S0158}
\end{qa}

\iin{北京大学}
\ith{时间线}
\ist{认知与规划}
\begin{qa}{M0048}{我们学校没有推免资格，我还能保研吗}
\textbf{本校没有推免资格，你就不能走推免。推免资格是学校层面的资质，不是个人能申请的。但你可以正常报考全国硕士研究生招生考试。}

能不能开展推免工作由教育部核定：高校认为本校符合条件（教学质量优秀；具有博士学位授予权，或具有硕士学位授予权且独立招收硕士研究生连续15年、体育艺术院校连续6年以上；招生工作秩序良好；办学行为规范）的，可向所在地省级研究生招生主管部门提出申请，省级主管部门核实遴选后转报教育部，教育部组织专家委员会评议后向社会公示，公示无异议的由教育部通知并公布开展推荐工作的高校名单（教学〔2006〕14号第十条、第十一条）。第八条也写明教育部公布可开展推荐工作的学校名单和年度推荐名额。

名单是动态的：2025年教育部办公厅印发《关于对新增推免资格高校予以备案的通知》，对北京电子科技学院等\textbf{67所}新增推免资格高校予以备案，其中有不少地方院校。
\nref{S0454}
\nref{S0640}
\nref{S0050}
\end{qa}

\iin{教育部学生服务与素质发展中心（学信网/研招网）}
\ith{资格认知}
\ist{认知与规划}
\begin{qa}{M0049}{独立学院有保研名额吗}
\textbf{不能。独立学院学生被全国统一文件明确排除在推免推荐对象之外。}

教学〔2006〕14号第十四条第（一）项把推荐对象限定为纳入国家普通本科招生计划录取的应届毕业生，括号里明确写\textbf{不含专升本、第二学士学位、独立学院学生}。这是全国统一规定，不是某一所学校的口径。

院校文件后来也反复重述，例如广东外语外贸大学2026年推荐优秀应届本科毕业生免试攻读硕士学位研究生实施办法，第二条在界定适用范围时明确不含独立学院学生。

要区分两件事：

\begin{itemize}
  \item \textbf{推免名额}由教育部下达给具有推免资格的推荐高校，独立学院不在推荐主体范围内，因此不存在独立学院自己的推免名额。
  \item \textbf{考研}不受影响：取得本科毕业证书和学士学位后，可按普通应届本科毕业生身份报名统考，录取后与推免生在同一培养体系内。
\end{itemize}

个别独立学院转设或并入母体高校后，学生身份按转设后的学校性质认定，需看转设批文。
\nref{S0454}
\nref{S0134}
\end{qa}

\iin{教育部学生服务与素质发展中心（学信网/研招网）}
\ith{资格认知}
\ist{认知与规划}
\begin{qa}{M0060}{在职的人还能走保研这条路吗}
\textbf{不能。推免的对象是全日制应届本科毕业生，在职人员不属于这个范围；在职读研走的是统考（含非全日制）这条线。}

教学〔2006〕14号第十四条第（一）项界定推荐对象为纳入国家普通本科招生计划录取的应届毕业生；同办法第三条对推荐的定义是普通高等学校按规定对本校\textbf{优秀应届本科毕业生}进行遴选。教育部历年硕士研究生招生工作管理规定中，推免生的资格审核与备案都围绕应届本科毕业生展开。在职人员已取得本科学历并脱离应届身份，无法进入推免通道。

还有一类容易混淆的免试通道：退役人员中，服役期间获得三等战功、二等功以上奖励或者二级以上表彰、符合全国硕士研究生招生考试报考条件的，可申请免初试攻读硕士研究生，须在教育部规定的全国统考报名时间内登录推免服务系统报名。这属于立功表彰退役军人免试生，与在职人员无关。

唯一的在职加免试特例是公费师范毕业生申请免试攻读非全日制教育硕士，那是履约培养项目，不是普通推免。
\nref{S0454}
\nref{S0389}
\nref{S0609}
\end{qa}

\iin{教育部学生服务与素质发展中心（学信网/研招网）}
\ith{在职}
\ist{认知与规划}
\section{信息渠道}
\begin{qa}{M0005}{保研信息一般从哪看？光看学院官网够吗}
\textbf{学院官网远远不够。至少要看四类渠道：教育部推免服务系统、目标院校研究生院与院系官网、省级教育考试机构公告、本校教务处与学院。}

各渠道各管一段：

\begin{enumerate}
  \item \textbf{教育部推免服务系统}（yz.chsi.com.cn/tm）：教学厅〔2014〕5号规定，推免生可通过该系统查询招生单位的招收推免生章程和专业目录、填写报考志愿、接收并确认复试及待录取通知，最终名单以系统备案信息为准。这是唯一有录取效力的平台。
  \item \textbf{目标院校研究生院与院系官网}：接收条件、材料清单、考核方式、复试时间都在这里。东北师范大学2027年办法直接要求考生密切关注中国研究生招生信息网、我校研究生院网站、各招生院部网站发布的通知公告。
  \item \textbf{省级教育考试机构公告}：统一的注册、志愿填报、确认时间都在这类公告里。四川省教育考试院2027年公告给出9月18日注册、9月21日9:00开始填报志愿、10月18日截止、10月20日前未落实接收单位不再保留资格等全部节点。
  \item \textbf{本校教务处与学院}：推免资格、排名口径、名额分配出自这里，校内资格申请往往更早，北京大学2027届为9月1日至7日。
\end{enumerate}
\nref{S0048}
\nref{S0607}
\nref{S0389}
\nref{S0553}
\end{qa}

\iin{中华人民共和国教育部}
\ith{信息渠道}
\ist{认知与规划}
\begin{qa}{M0008}{双非真的能保到985吗？我们学院去年好像只有一个去了211}
\textbf{能，但难度取决于目标院校接收端的门槛。政策上没有任何条款禁止双非学生外推，真正的门槛在院系申请条件里。}

政策面：教学〔2006〕14号第七条提倡优势互补、鼓励推免生在不同地区间、招生单位间及不同学科间交流；教学厅〔2014〕5号规定所有推免生均享有依据招生政策自主选择报考招生单位和专业的权利，所有推免名额（除有特殊政策要求的专项计划外）均可向其他招生单位推荐，推荐高校不得以任何形式限制推免生自主报考。本科出身不构成规则障碍。

门槛在接收端。山东大学控制科学与工程学院2026年夏令营通知按学校层次给了不同排名要求：

\par\vspace{2pt}\nobreak\begin{center}\footnotesize
\begin{tabularx}{\textwidth}{>{\raggedright\arraybackslash}X>{\raggedright\arraybackslash}X}
\toprule
\textbf{申请者学校类别} & \textbf{成绩排名要求} \\ \midrule
985工程高校 & 专业排名前20\%，并有望获得推免资格 \\
211工程高校 & 专业排名前10\%，并有望获得推免资格 \\
其他具有研究生推免资格高校 & 专业排名前8\%，并有望获得推免资格 \\
\bottomrule
\end{tabularx}
\end{center}\vspace{2pt}\par

双非要拿到同一张入场券，需要更高的校内排名。务实做法是把目标院校分成冲、稳、保三档，重点投往年确实录取过你本科学校学生的院系。
\nref{S0454}
\nref{S0048}
\nref{S0102}
\end{qa}

\iin{教育部学生服务与素质发展中心（学信网/研招网）}
\ith{双非认知}
\ist{认知与规划}
\begin{qa}{M0020}{保研有没有卷不卷的说法，怎么判断我们学院卷不卷}
\textbf{卷不卷没有官方定义，但可用三个可查的客观量判断：本院名额数、往年公示名单规模、同学的报名意愿。}

\begin{enumerate}
  \item \textbf{名额数}：本院当年推免细则或学院通知里的名额分配。湖南大学2027年通知写明学校按教育部下达的推免生名额分配至当年度有应届本科毕业生的学院，学院再往下切。名额和应届毕业生人数一比，基本盘就出来了。
  \item \textbf{往年名单规模}：各校教务处和学院会公示拟推荐名单、递补名单。浙江师范大学2027年公示的拟推荐名单、递补名单能看到具体人数和专业分布；中南大学材料科学与工程学院也公示过预推荐名单和增补名单。连续看三年能看出名额是稳定还是被压缩。
  \item \textbf{报名意愿}：只能靠问，但有间接信号，例如成都大学商学院细则规定候补人选按各专业名额数的50\%计算，候补比例高说明边缘选手多、竞争激烈。
\end{enumerate}

另一个易被忽略的维度是外推意愿：如果本院历年名单里大量学生最终留本校，说明外推渠道不畅，你的卷其实是在一个更小的池子里。
\nref{S0540}
\nref{S0571}
\nref{S0275}
\nref{S0079}
\end{qa}

\iin{湖南大学}
\ith{卷度判断}
\ist{认知与规划}
\begin{qa}{M0021}{保研失败的话还有退路吗？能不能转头考研}
\textbf{有退路。推免失败（未落实接收单位）不影响本科学籍，并且可在统考报名截止前转报考当年研究生招生考试。}

两点明确的官方口径：

\begin{itemize}
  \item 中国政府网转北京教育考试院问答《推免生还能参加统考吗？》：在规定截止日期未被招生单位录取的推免生不再保留推免资格，可在网上报名截止时间内选择报考当年全国硕士研究生招生考试。已被招生单位录取的推免生，不得再报名参加当年全国硕士研究生招生考试。
  \item 四川省教育考试院2027年公告：2026年10月20日前未落实接收单位的推免生不再保留推免生资格，但可报名参加全国硕士研究生招生考试。
\end{itemize}

也就是说，\textbf{只有已确认待录取的推免生才被锁死}；未落实接收单位的，10月下旬之后仍可转入统考。

时间差很关键：推免结果落定在10月20日前后，而统考网上报名通常也在10月进行，窗口贴得很近。排名在边缘时，提前备好统考复习资料与报名材料更稳妥。
\nref{S0352}
\nref{S0389}
\end{qa}

\iin{中国政府网}
\ith{风险预案}
\ist{认知与规划}
\begin{qa}{M0029}{医学生保研和临床实习会冲突吗？怎么安排时间}
\textbf{现行公开文件中未见全国统一的医学推免生实习安排规定，冲突的解决方式由各医学院校自行安排，必须以本校当年文件为准。}

从可查到的官方文件看，医学类院校的推免通知主要规定接收条件、专业要求与考核方式，不涉及临床实习与推免时间的统筹。例如南昌大学基础医学院2027年接收推免生通知列出的申请条件是专业要求（医学类或生物类相关专业）、体检要求与材料要求；复旦大学临床医学院2026年推免生招生预报名通知同样只规定报名与考核安排；深圳大学医学部2027级推免研究生公告亦属接收端文件。

制度层面能确定的是推免的硬时间窗：2027届志愿填报在2026年9月21日至10月18日，待录取确认从9月24日开始（四川省教育考试院2027年公告）；夏令营集中在6—7月。实习期与之重叠时，需要与学院教务部门协调请假或调整轮转。

各校附属医院轮转制度差异极大，本条不给通用建议。
\nref{S0001}
\nref{S0172}
\nref{S0265}
\nref{S0389}
\end{qa}

\iin{南昌大学江西医学院基础医学院}
\ith{医学学制}
\ist{认知与规划}
\begin{qa}{M0036}{我们学院保研细则每年都改，以哪一版为准}
\textbf{以学校当年、面向你这一届公布的推免工作实施办法和学院工作细则为准。细则每年修订是常态，不存在通用版。}

两条制度支撑：

\begin{itemize}
  \item 教学〔2006〕14号第十五条第（一）项要求学校按照本办法的原则和规定，制订\textbf{年度}推荐工作实施办法，于推荐工作启动前在校内各院（系）公布。
  \item 教学厅〔2014〕5号要求学校和院系推免生遴选办法均要严格遵守相关文件和推免政策管理规定，广泛征求师生意见，与推免名额一并提前公布并严格执行。
\end{itemize}

院校也是这样运作的：湖南大学每年由教务处发布《关于公布各学院推荐当年优秀应届本科毕业生免试攻读硕士学位研究生工作细则的通知》，逐年更新；北京大学2027届办法则写明如上级部门出台新的推免政策或工作要求，本办法将做相应调整并另行通知。

判断版本的顺序：先确认标题年份是不是你毕业那一届，再看名额、时间节点、公示安排是不是本年度，最后核对落款单位与日期。
\nref{S0454}
\nref{S0048}
\nref{S0539}
\nref{S0553}
\end{qa}

\iin{教育部学生服务与素质发展中心（学信网/研招网）}
\ith{细则变动}
\ist{认知与规划}
\begin{qa}{M0045}{保研后能不能反悔去工作}
\textbf{能反悔，但代价自负，而且有几条明确的锁死条款：一旦在推免服务系统确认待录取，你当年就不能再报名统考。}

四条约束：

\begin{enumerate}
  \item \textbf{全国统一}：被确定的推免生如果无招生单位接收或本人放弃，则该生推免生资格作废，名额不得转让（教学〔2006〕14号第十三条）。
  \item \textbf{系统层面}：接受待录取通知后不能再填报其他志愿或接受其他复试、待录取通知，其他招生单位也不能再向该推免生发放复试或待录取通知；未通过系统确认接受待录取通知则录取无效（四川省教育考试院2027年公告、上海财经大学2027年操作说明）。确需取消的，须遵守推免政策规定向招生单位提出申请，招生单位同意后可通过系统取消。
  \item \textbf{院校层面}：北京大学2027届办法写明已落实接收单位的推免生不得再申请出国（境）留学、参加就业派遣等，否则违约责任自负，教务部也不为这类学生出具出国成绩单及学历证明；东北师范大学2027年办法写明如因个人原因放弃录取资格，我校不予出具任何证明和手续。
  \item \textbf{时间成本}：推免结果落定在10月中下旬，此时统考报名已近尾声。
\end{enumerate}
\nref{S0454}
\nref{S0389}
\nref{S0182}
\nref{S0553}
\nref{S0607}
\end{qa}

\iin{教育部学生服务与素质发展中心（学信网/研招网）}
\ith{违约}
\ist{认知与规划}
\begin{qa}{M0051}{保研拿到的offer有法律效力吗？学校能反悔吗}
\textbf{真正生效的是推免服务系统里的确认动作，不是夏令营或预推免的口头意向；学校也不能随意反悔，取消录取必须落回明文规定的情形。}

推免生一侧的效力规则：

\begin{itemize}
  \item 招生单位通过系统发送待录取通知，推免生须在规定时间内答复；未按时回复的，招生单位可取消待录取通知。
  \item \textbf{未通过系统确认接受待录取通知，则录取无效}；确认接受后不能再填报其他志愿或接受其他复试、待录取通知。
  \item 学校层面的拟录取也要公示：招生单位对推免生申请者进行复试，确定拟录取名单并公示7天以上，无异议的由招生单位通知该考生（教学〔2006〕14号第十八条）。
\end{itemize}

学校可合法取消的情形是写明的。浙江大学2027年通知列出五项：思想政治素质和品德考核不合格或政审不合格；从推免生报名至入学报到之日受到任何处分；本科最后一学年计入毕业要求学分的课程出现不及格记录；本科毕业设计（论文）成绩未达到良好及以上（或80分及以上）；规定的新生入学日期前无法提供毕业证书和学士学位证书。教学〔2006〕14号第二十条也规定，入学前未取得学士学位或本科毕业证书，或受处分的，取消录取资格。
\nref{S0389}
\nref{S0182}
\nref{S0454}
\nref{S0032}
\end{qa}

\iin{www.sceea.cn}
\ith{违约}
\ist{认知与规划}
\begin{qa}{M0052}{保研信息差真的那么大吗？怎么才能不落后}
\textbf{信息差确实存在，但它的来源不是有秘密文件，而是四类公开信息你有没有系统地看：省级统一时间表、目标院校院系通知、本校细则、系统里的公开数据。}

制度上的要求是信息公开，也就是说差距在于看没看到：

\begin{itemize}
  \item 教学厅〔2014〕5号要求做好推免名额、推荐办法、推免生名单（含姓名、院系、综合测评成绩等）、复试录取办法、拟录取推免生名单（含姓名、复试成绩等）、咨询申诉渠道等推免招生重要信息的公示公开工作。
  \item 教学〔2006〕14号第二十二条要求推荐及接收单位公开推免生政策规定、推免生资格、招收名额、录取信息、考生咨询及申诉渠道。
  \item 推免生还可在规定时限内查看本人报考并参加复试的招生单位同专业、同层次、同类别推免生录取公示信息（四川省教育考试院2027年公告）。
\end{itemize}

容易被漏掉的两块：省级教育考试机构的公告，以及目标院校的预推免系统——很多学校在教育部系统开放前就有自己的报名系统，如湖南大学2027年预报名8月28日起。
\nref{S0048}
\nref{S0454}
\nref{S0389}
\nref{S0158}\par\textbf{本版实操补充：} 把“信息差”拆成可核实的差异：目标单位当年接收办法、资格门槛、申请窗口、考核形式和名额变化。每周固定一次查目标院系官网与研招网，把文件下载后按“年度\_单位\_文件类型”命名，并在台账记发布日期、适用对象和原始链接。群聊转述只用作线索，找不到原文就标成待核实，不据此删掉学校或支付服务费。
\end{qa}

\iin{中华人民共和国教育部}
\ith{信息渠道}
\ist{认知与规划}
\section{推免系统}
\begin{qa}{M0001}{保研到底是怎么个流程？是不是大四开学直接填个志愿就行了}
\textbf{不是大四开学填个志愿就行。填志愿之前，你必须先在校内拿到推免资格；这是两个系统、两个环节。}

以2027届为例，完整链条是：

\begin{enumerate}
  \item \textbf{校内资格申请（大四开学初）}：北京大学2027届为2026年9月1日至7日开通校内推免资格申请系统，9月8日至10日院系初审，9月11日报送，9月18日前学校审核并公示，9月19日起学生可查结果。这一步决定你有没有推免生身份。
  \item \textbf{注册推免服务系统}：2026年9月18日起注册、填报个人资料并查看本人推免资格。
  \item \textbf{填报志愿}：9月21日9:00开始填报，10月18日截止。可同时填报三个平行志愿，每个志愿提交后48小时内不允许修改。
  \item \textbf{确认复试通知}：9月22日8:00开始确认招生单位发送的复试通知；未通过系统确认接受复试通知的考生不能被录取。
  \item \textbf{确认待录取}：9月24日14:00开始确认待录取通知。接受后不能再填报其他志愿或接受其他复试、待录取通知。
  \item \textbf{收尾}：10月20日前未落实接收单位的推免生不再保留推免资格，但可报名参加全国硕士研究生招生考试。
\end{enumerate}

整条链上，志愿和确认都必须在\textbf{教育部推免服务系统}里完成：未经过该系统接收的推免生一律无效。
\nref{S0553}
\nref{S0389}
\nref{S0607}
\end{qa}

\iin{北京大学}
\ith{保研认知}
\ist{认知与规划}
\begin{qa}{M0006}{别人都在说夏令营预推免九推，这三个到底啥区别}
\textbf{三者是推免接收在时间上先后不同的三个阶段，共同点是都必须回到教育部推免服务系统才算数。}

\par\vspace{2pt}\nobreak\begin{center}\footnotesize
\begin{tabularx}{\textwidth}{>{\raggedright\arraybackslash}X>{\raggedright\arraybackslash}X>{\raggedright\arraybackslash}X>{\raggedright\arraybackslash}X}
\toprule
\textbf{阶段} & \textbf{时间（2027届）} & \textbf{谁组织} & \textbf{效力} \\ \midrule
夏令营 & 6—7月举办，报名截止多在6月中下旬 & 招生单位院系自主 & 各校不同，优秀营员有的对接待录取，有的只是免筛选 \\
预推免 & 8—9月，系统开放前 & 招生单位院系自主，多用本校报名系统 & 院系意向，仍须回系统走流程 \\
九推（正式填报） & 9月21日9:00开始填报志愿，10月18日截止 & 教育部推免服务系统 & 唯一有效的报考录取渠道 \\
\bottomrule
\end{tabularx}
\end{center}\vspace{2pt}\par

制度上的关键一句是：未经过教育部推免服务系统接收的推免生一律无效（东北师范大学2027年办法）。所以夏令营和预推免拿到的都是意向，最终都要在系统里完成填报志愿、接受复试通知、接受待录取三步。例如湖南大学2027年预报名8月28日起在自己的报名信息采集系统里进行，但申请人仍须在全国推荐免试攻读研究生信息公开管理服务系统中正式报名。

简单记：\textbf{夏令营是提前认识，预推免是提前锁定，九推是正式成交}。只有最后一步在系统里完成，录取才生效。
\nref{S0389}
\nref{S0607}
\nref{S0158}
\nref{S0264}
\end{qa}

\iin{www.sceea.cn}
\ith{阶段认知}
\ist{认知与规划}
\begin{qa}{M0009}{保研名额是国家给的还是学校自己分的？为什么我们学院今年多了一个}
\textbf{名额总量由国家控制、按校下达，学校再分配到你所在的学院和专业。学院今年多了一个，来自学校内部的二次分配。}

全国层面：教育部负责全国推免工作的宏观管理，规定开展推荐工作的学校标准和条件，公布可开展推荐工作的学校名单和年度推荐名额；教育部对年度推免生规模实行总量控制（教学〔2006〕14号第八条、第九条）。名额确定原则按学校层次给出比例（第十二条）：设研究生院的高校一般按应届本科毕业生数的15\%左右、未设研究生院的211工程高校约5\%、其他高校约2\%，初次开展的前3年约1\%；人文、理科等人才培养基地班按基地班招生人数的50\%左右单独增加。

学校内部的分配依据由学校自定：湖南大学2027年通知写明学校按教育部下达名额分配至有应届本科毕业生的学院，研究生支教团等国家专项计划名额按国家相应政策分配至相关推荐单位；广东外语外贸大学2026年办法则综合考虑各学院应届毕业生人数、硕士点、本科专业数量以及上年度名额使用情况。年际变动通常来自毕业生人数变化、上年度使用情况和是否新增专项计划。
\nref{S0454}
\nref{S0540}
\nref{S0134}
\end{qa}

\iin{教育部学生服务与素质发展中心（学信网/研招网）}
\ith{名额来源}
\ist{认知与规划}
\begin{qa}{M0011}{学硕和专硕在保研里有什么区别？听说专硕好保一点}
\textbf{从2014年起，推免名额在下达时不再区分学术学位和专业学位，也不再有留校限额。但学硕与专硕在接收端仍是两条不同轨道。}

全国统一的规则变化：教学厅〔2014〕5号明确，2014年起教育部下达推免名额时不再区分学术学位和专业学位，不再设置留校限额；推荐高校也不得对本校推免名额限制学术学位与专业学位报考类型，不得自行设置留校限额或名额。也就是说在\textbf{推荐端}不存在学硕名额和专硕名额之分，你能不能被推荐只看校内排名。

差别出现在\textbf{接收端}：

\begin{itemize}
  \item 招生单位的专业目录会分别列出学术学位与专业学位专业的接收人数，考核内容、学制、学费都可能不同。东北师范大学2027年办法里直博生基本学制5年、学费每学年1万元，学制与收费另见学校文件。
  \item 中山大学2027年推免办法写明全日制硕士研究生学制一般为2—3年，推免生入学时均须缴纳学费。
\end{itemize}

专硕好保一点的说法通常来自专业报名冷热不均，而不是制度上给专硕留了更容易的名额。
\nref{S0048}
\nref{S0607}
\nref{S0167}
\end{qa}

\iin{中华人民共和国教育部}
\ith{学硕专硕}
\ist{认知与规划}
\begin{qa}{M0028}{港澳台学生保研是走什么渠道}
\textbf{现行公开文件中未见全国统一的港澳台学生推免专门规定。港澳台侨学生若在内地高校全日制就读，能否推免取决于该校细则是否列入推荐对象，必须以本校当年文件为准。}

可以查到两种口径：

\textbf{明确包含}：成都大学商学院2027届推免遴选细则把推荐对象写为本校通过普通高考（含保送方式）招生录取或面向港澳台华侨学生招生录取的全日制应届本科毕业生；武汉轻工大学2027年接收推免生工作办法的申请条件同样表述为应届本科毕业生，通过普通高考（含保送方式）招生录取或面向港澳台华侨学生招生录取的全日制本科生。

\textbf{未提及}：北京大学2027届办法的申请条件只列纳入国家普通本科招生计划录取的应届毕业生（不含专升本、第二学士学位学生），未单独提及港澳台学生；教学〔2006〕14号第十四条也只写纳入国家普通本科招生计划录取的应届毕业生。

读研渠道要区分：港澳台学生另有面向港澳台地区的研究生招生考试，与推免是两套体系；教育部令第41号《普通高等学校学生管理规定》第六十六条明确对港澳台侨学生、留学生的管理参照本规定执行，而非直接适用。
\nref{S0079}
\nref{S0409}
\nref{S0553}
\nref{S0454}
\nref{S0040}
\end{qa}

\iin{成都大学商学院}
\ith{港澳台}
\ist{认知与规划}
\begin{qa}{M0037}{保研是不是有个推免系统？什么时候开放}
\textbf{有，就是教育部推免服务系统（全国推荐免试攻读研究生信息公开管理服务系统），网址 yz.chsi.com.cn/tm。2027届9月18日起注册并查看资格，9月21日9:00起填报志愿。}

具体时间（以2027届为例，四川省教育考试院2026年9月21日公告）：

\begin{itemize}
  \item \textbf{2026年9月18日起}：注册、填报个人资料信息并查看本人推免资格。
  \item \textbf{2026年9月21日9:00}：开始填报志愿。
  \item \textbf{10月18日}：截止填报志愿。
  \item \textbf{9月22日8:00}：开始确认招生单位发送的复试通知。
  \item \textbf{9月24日14:00}：开始确认招生单位发送的待录取通知。
  \item \textbf{10月20日}：未落实接收单位的推免生不再保留推免生资格。
\end{itemize}

上海财经大学2027年操作说明也表述为9月21日至10月20日，推免生可通过推免服务系统填写报考志愿，接收并确认招生单位的复试及待录取通知。

这个系统的地位是排他的：教学厅〔2014〕5号规定，推免生（含推免硕士生和直博生）资格审核确认、报考、录取以及备案公开等相关工作均须通过该系统进行，最终推免生名单以系统备案信息为准；东北师范大学2027年办法也写明未经过教育部推免服务系统接收的推免生一律无效。开通时间每年略有调整，以教育部当年公告和本省公告为准。
\nref{S0389}
\nref{S0182}
\nref{S0048}
\nref{S0607}
\end{qa}

\iin{www.sceea.cn}
\ith{推免系统}
\ist{认知与规划}
\begin{qa}{M0038}{推免服务系统里可以填几个志愿}
\textbf{可以同时填报三个平行志愿，不区分先后顺序；每个志愿提交后48小时内不允许修改。通过研究生支教团专项计划遴选的推免生只能报考本校专业志愿。}

依据四川省教育考试院2027年公告：推免生（有特殊政策要求的专项计划推免生除外）可同时在系统填报3个平行志愿（不区分先后顺序），每个志愿在提交后的48小时内不允许修改。若招生单位在48小时内通过系统拒绝申请，或48小时后推免生仍未接到复试通知、复试未通过，则该志愿自动解锁，可继续填报其他志愿。

上海财经大学2027年操作说明补充：若推免生拒绝招生单位发送的待录取通知，该志愿自动解锁，可继续填报其他志愿；通过支教团专项计划遴选的推免生\textbf{只能选择报考本校专业志愿}。

志愿数量不等于录取数量：推免生只能接受\textbf{一个}待录取通知，确认接受后不能再填报其他志愿或接受其他复试、待录取通知，未通过系统确认接受待录取通知则录取无效。所以三个志愿应按冲、稳、保分层投放，并留意48小时锁定期占用窗口时间。
\nref{S0389}
\nref{S0182}
\nref{S0444}
\end{qa}

\iin{www.sceea.cn}
\ith{推免系统}
\ist{认知与规划}
\begin{qa}{M0039}{保研一定要参加预推免吗？只在夏令营拿到offer行不行}
\textbf{不一定。制度上没有必须参加预推免的环节，最终生效的只有教育部推免服务系统里的一整套操作；但如果目标院校要求预推免，你就得参加。}

三件事要区分：

\begin{enumerate}
  \item \textbf{夏令营与预推免是招生单位的自主行为}：东北师范大学2027年办法写明，考生是否必须先在本校系统报名、以及院部如何组织考核由各招生院部自定；同时未在学校系统中报名的申请者如有意愿报考，可在教育部推免服务系统中报名，报考院部将结合实际确定是否再次组织综合考核。这说明预推免不是必经环节。
  \item \textbf{推免服务系统才是生效环节}：同一办法明确，已在学校系统中报名的申请者仍须在教育部推免服务系统中报名，\textbf{未经过该系统接收的推免生一律无效}。
  \item \textbf{志愿和确认有明确时间窗}：2027届9月21日9:00开始填报志愿、10月18日截止填报、9月22日8:00起确认复试通知、9月24日14:00起确认待录取通知（四川省教育考试院2027年公告）。
\end{enumerate}

所以只在夏令营拿到意向行不行，答案是：可以，前提是对方意向能衔接上系统里的待录取流程，且你按要求完成填报与确认。风险在于对方当年只通过预推免批次锁人，那就可能错过窗口。
\nref{S0607}
\nref{S0389}
\nref{S0182}
\end{qa}

\iin{东北师范大学}
\ith{路径选择}
\ist{认知与规划}
\begin{qa}{M0059}{我在国外读的本科，能申请国内高校的推免吗}
\textbf{不能。推免的对象是境内高校纳入国家普通本科招生计划录取的应届毕业生，在国外读本科不属于这个范围。}

教学〔2006〕14号第十四条第（一）项界定的推荐对象是纳入国家普通本科招生计划录取的应届毕业生；推荐行为由具有推免资格的高等学校对本校学生组织实施（同办法第三条：推荐是指普通高等学校按规定对本校优秀应届本科毕业生进行遴选）。国外高校不在教育部公布的开展推免工作的高校名单内，也就不存在推荐主体。

国外本科毕业生读国内研究生的路径是另一套：

\begin{itemize}
  \item 参加全国硕士研究生招生考试。招生简章通常要求在境外获得学历学位的考生提交教育部留学服务中心出具的《国（境）外学历学位认证书》，例如厦门大学2026年复试录取工作办法、福建农林大学2026年硕士招生简章都有此类要求。
  \item 或者申请部分院校的申请考核制、国际化项目等，以目标院校当年招生简章为准。
\end{itemize}

另外，教学〔2006〕14号第二十条规定，入学前未取得学士学位或本科毕业证书，或受处分的，取消录取资格。境外学历认证的时间要预留出来。
\nref{S0454}
\nref{S0413}
\end{qa}

\iin{教育部学生服务与素质发展中心（学信网/研招网）}
\ith{留学生}
\ist{认知与规划}
\section{保内保外}
\begin{qa}{M0007}{保研是不是只能保本校？想出去是不是很难}
\textbf{不是只能保本校。全国统一规定明确保障推免生自主选择报考招生单位和专业的权利，推荐高校不得以任何形式限制；留校限额从2014年起就取消了。}

\begin{itemize}
  \item 教学厅〔2014〕5号：所有推免生均享有依据招生政策自主选择报考招生单位和专业的权利，所有推免名额（除有特殊政策要求的专项计划外），均可向其他招生单位推荐。推荐高校要充分尊重并维护考生自主选择志愿的权利，不得将报考本校作为遴选推免生的条件，也不得以任何其他形式限制推免生自主报考。
  \item 同文件：2014年起，教育部下达推免名额时不再区分学术学位和专业学位，不再设置留校限额；推荐高校也不得对本校推免名额限制学术学位与专业学位报考类型，不得自行设置留校限额或名额。
\end{itemize}

这里有一处新旧关系要说明：教学〔2006〕14号第十七条曾规定设有研究生院的高校接收本校推免生的人数不得超过本校推免生总数的65\%（西部省份或军工、矿业、石油、地质、农林等特殊类型高校可放宽至75\%）。这个留校比例上限已被2014年通知的不再设置留校限额取代，引用时应以2014年文件为准。

所以出去难不难，卡点不在政策，而在目标院校的接收条件和你自己的准备。
\nref{S0048}
\nref{S0454}
\end{qa}

\iin{中华人民共和国教育部}
\ith{保内保外}
\ist{认知与规划}
\begin{qa}{M0014}{保研和出国同时搞现实吗？语言考试和夏令营时间撞了怎么办}
\textbf{可以同时准备，但有一处硬冲突：多所高校把不计划本科毕业后赴境外留学写进了推免申请条件，并且已落实接收单位的推免生会被限制出国。}

北京大学2027届推免办法的相关条款最具体：

\begin{itemize}
  \item 申请条件之一为不计划本科毕业后赴境外留学；
  \item 特别强调已计划毕业后出国（境）留学的学生，不能申请推免资格；
  \item 已落实接收单位的推免生，不得再申请出国（境）留学、参加就业派遣等，否则违约责任自负，教务部也不为这类学生出具出国（境）申请所需的成绩单及学历证明。
\end{itemize}

时间上的冲突：夏令营集中在6—7月，如浙江大学传媒与国际文化学院2026年夏令营7月7—8日、华中科技大学网络空间安全学院7月9—10日，而语言考试和申请材料准备通常在同期进行，精力必然互相挤占。

所以务实判断是：把两条路当成互斥选项来规划。以出国为主就不要再申请推免资格占用名额；以推免为主，出国准备应放在推免结果落定（10月下旬）之后。
\nref{S0553}
\nref{S0003}
\nref{S0013}
\end{qa}

\iin{北京大学}
\ith{出国对比}
\ist{认知与规划}
\begin{qa}{M0043}{保研去外校，本校不放人怎么办}
\textbf{如果本校以外推为条件卡你，那是违规的。全国统一规定明确禁止把报考本校作为遴选条件，也禁止设置留校限额；拿规则说话并走申诉渠道。}

可援引的条款：

\begin{itemize}
  \item 教学厅〔2014〕5号：所有推免生均享有依据招生政策自主选择报考招生单位和专业的权利，所有推免名额（除有特殊政策要求的专项计划外）均可向其他招生单位推荐；推荐高校不得将报考本校作为遴选推免生的条件，也不得以任何其他形式限制推免生自主报考。
  \item 同文件：推荐高校不得对本校推免名额限制学术学位与专业学位报考类型，不得自行设置留校限额或名额。
  \item 教学〔2006〕14号第二十三条：高等学校应当将推免生工作中学生的申诉，纳入校内申诉渠道。
\end{itemize}

申诉路径通常是：先向学院推免工作小组提交书面异议（公示名单里一般会留联系人和电话），再向学校教务处或推免生遴选工作领导小组反映，必要时向学校所在省级教育行政部门或省级教育招生考试机构申诉。四川省2027年公告明确，考生如有疑问可向推荐高校、招生单位所在省级教育行政部门、省级教育招生考试机构申诉。

前提是你的推荐资格本身站得住，排名、材料、公示流程都合规，申诉才有力。
\nref{S0048}
\nref{S0454}
\nref{S0389}
\end{qa}

\iin{中华人民共和国教育部}
\ith{保内保外}
\ist{认知与规划}
\begin{qa}{M0046}{保研失败是被退学吗？还是只是没名额}
\textbf{不是退学。推免失败只意味着你没有拿到推免资格或没有落实接收单位，本科学籍和毕业完全不受影响。}

可以确认的后果只有一条：不再保留推免资格。中国政府网转北京教育考试院的问答写明，在规定截止日期未被招生单位录取的推免生不再保留推免资格，可在网上报名截止时间内选择报考当年全国硕士研究生招生考试；四川省教育考试院2027年公告也写明，2026年10月20日前未落实接收单位的推免生不再保留推免生资格，但可报名参加全国硕士研究生招生考试。

推免资格作废的全国统一规则出自教学〔2006〕14号第十三条：截至教育部规定时间后，被确定的推免生如果无招生单位接收或本人放弃，则该生推免生资格作废，名额不得转让。

真正会失去东西的情形是另外几种，且与推免失败无关：受到处分、学术不端、入学前未取得学士学位或本科毕业证书。教学〔2006〕14号第二十条规定，这些情况下取消的是\textbf{录取资格}，而不是学籍。
\nref{S0352}
\nref{S0389}
\nref{S0454}
\end{qa}

\iin{中国政府网}
\ith{后果认知}
\ist{认知与规划}
\begin{qa}{M0056}{保研里说的全聚德是什么意思？被全拒了怎么办}
\textbf{全聚德是学生圈里的戏称，官方文件中没有这个词，指投出去的夏令营、预推免和系统志愿全部被拒。制度上仍有兜底路径。}

先说明用词：教学〔2006〕14号、教学厅〔2014〕5号等全国性文件，以及各省公告、各校办法中均无全聚德一词，它是民间对全被拒的谐音说法。写进手册是为了让你看懂大家在说什么，不是规则。

真被全拒之后，规则上还剩的空间：

\begin{itemize}
  \item \textbf{志愿会解锁、可以继续填}：填报的志愿提交后48小时内不允许修改；若招生单位在48小时内拒绝，或48小时后仍未接到复试通知、复试未通过，则该志愿自动解锁，可继续填报其他志愿（四川省教育考试院2027年公告）。
  \item \textbf{填报窗口到10月18日}：2027届志愿填报从9月21日9:00到10月18日截止，期间可以多轮投递。
  \item \textbf{10月20日是最后节点}：未落实接收单位的推免生不再保留推免生资格，但可报名参加全国硕士研究生招生考试。
\end{itemize}

所以全拒之后的两条路是：继续投递仍有缺额的院校（通常是非热门院校和中西部院校），或者放弃推免转统考。两条路都需要在10月下旬前做出决定。
\nref{S0389}
\nref{S0182}
\nref{S0352}
\nref{S0454}
\nref{S0048}\par\textbf{本版实操补充：} 先把失败拆成三类复盘：资格或硬门槛不符、材料表达没有证明匹配度、面试考核失误。逐所记录拒信时间、招生方向和已投入环节，问招生办是否还有后续批次或候补安排；不要把未回复自动当作正式拒绝，也不要在未落实接收前停止关注其他路径。若本校资格仍在，立刻确认系统窗口和统考报名节点；是否转考研以当年教育部及省级公告为准。
\end{qa}

\iin{www.sceea.cn}
\ith{圈内用语}
\ist{认知与规划}
\section{学科差异}
\begin{qa}{M0017}{保研是不是一定要有科研？没进过实验室是不是没戏}
\textbf{全国规定没有必须有科研这条硬门槛，但把创新意识、创新能力和专业能力倾向写成了考查重点。没进过实验室不等于没戏，前提是学业指标和考核表现能补上。}

全国层面的表述是教学〔2006〕14号第六条：推免生工作应在对考生平时学习和科研能力综合测评基础上，突出对考生创新精神、创新能力和专业能力倾向等的考查；第十四条第（三）项要求推免生学术研究兴趣浓厚，有较强的创新意识、创新能力和专业能力倾向。同条第（六）项还留了口子：对有特殊学术专长或具有突出培养潜质者，经三名以上本校本专业教授联名推荐并经学校推免生遴选工作领导小组严格审查，可不受综合排名限制，但说明材料和推荐信要进行公示。

接收端常把科研当筛选条件。西安交通大学网络空间安全学院2026年夏令营把在科研、国家级科技竞赛等方面已取得突出成果列为一项条件，但同一份通知也给了排名路径：综合成绩排名在专业前一定比例同样可以入选。

所以科研是加分项而非唯一通道，真正的一票否决通常是不及格记录、处分和排名掉出范围。
\nref{S0454}
\nref{S0132}
\end{qa}

\iin{教育部学生服务与素质发展中心（学信网/研招网）}
\ith{科研认知}
\ist{认知与规划}
\begin{qa}{M0022}{保研要花很多钱吗？跑夏令营是不是特别烧钱}
\textbf{没有全国统一的花费标准，但有两笔确定成本：推免生需网上支付报名费，以及参加复试的差旅费用。是否烧钱，主要看投递范围和是否线下参加。}

制度上能查到的只有报名费：推免生在推免服务系统注册后应填报个人资料信息、上传个人照片并进行网上支付（四川省教育考试院2027年公告的推免生操作流程第（二）项）；2025年推免服务系统操作流程同样写明推免生需网上支付报名费。各校处理方式不同，上海财经大学2027年操作说明明确，教育部规定推免生应网上支付报名费，但根据我校安排本校推免生一律免收相关费用。

其余开销（跨省参加线下夏令营和复试的交通住宿、材料打印邮寄、语言考试）在公开文件中没有统一规定，属个人支出。北京大学力学与工程科学学院2026年夏令营通知甚至要求纸质材料用EMS寄送并明确签收时间，这类要求会产生实际成本。

降低成本靠地理集中投递，或优先选择线上举办的营，例如华中科技大学网络空间安全学院2026年夏令营明确线上举办。
\nref{S0389}
\nref{S0444}
\nref{S0182}
\nref{S0264}
\nref{S0003}
\end{qa}

\iin{www.sceea.cn}
\ith{机会成本}
\ist{认知与规划}
\begin{qa}{M0030}{文科生保研是不是比理工科难？没有实验室没有论文}
\textbf{制度上没有文科更难的条款，名额比例按全校应届本科毕业生数统一核定；文科的真实难点在评价维度更依赖绩点、论文和外语，缺少实验室这类可量化产出。}

制度层面能查到的差异只有一条：教学〔2006〕14号第十二条第（四）项规定，经教育部确定的人文、理科等人才培养基地的高等学校，按教育部批准的基地班招生人数的50\%左右单独增加推免生名额。这是对人文、理科基地班的\textbf{增加}，不是限制。除此之外，名额按学校层次和应届本科毕业生数核定，不因学科门类区别对待。

文科生实际面对的三个结构性难点：

\begin{enumerate}
  \item \textbf{成果形态不同}：理工科可用实验室经历、竞赛、专利证明能力，文科更依赖绩点排名、学术论文和写作样本，本科阶段更难产出。
  \item \textbf{外语门槛客观存在}：多校把六级写成硬条件。广东外语外贸大学2026年办法对非外语类专业要求全国大学英语六级不低于425分（或托福95分、雅思7.0分）。
  \item \textbf{考核更依赖表达}：接收端复试多为面试与专业笔试，对口头表达和知识面要求高。
\end{enumerate}
\nref{S0454}
\nref{S0134}
\end{qa}

\iin{教育部学生服务与素质发展中心（学信网/研招网）}
\ith{学科差异}
\ist{认知与规划}
\begin{qa}{M0031}{经管的保研是不是特别卷？要不要考个CPA加分}
\textbf{经管类竞争激烈是报名规模决定的，不是制度设了更高的门槛。至于CPA，现行公开的推免文件中未见列为通用加分项，是否算加分要看本校细则里竞赛奖项、职业资格的认定范围。}

先看制度：教学〔2006〕14号第六条要求在对平时学习和科研能力综合测评的基础上择优选拔；北京大学2027届办法允许把高水平科研创新成果、竞赛奖项等纳入综合评价体系，但强调要按教育部要求严格审核认定并对相关材料存档，且不得单列名额、单独排序；广东外语外贸大学2026年办法把本科阶段学业综合成绩规定为推免工作最基础的遴选指标。

这意味着任何加分项都必须与学业成绩一起排序，且加分指标的认定范围、标准要提前公布。云南师范大学2027年通知要求各单位组织专家论证并提出本学科高水平科研创新成果、竞赛奖项等加分指标的认定范围、标准，并提前公布。

CPA属于职业资格考试，不是学科竞赛或科研成果，在公开细则中很少被列入加分目录。它起作用的场合通常是商学院在复试中评价专业能力时作为佐证材料，而不是直接加分。
\nref{S0454}
\nref{S0553}
\nref{S0134}
\nref{S0568}
\end{qa}

\iin{教育部学生服务与素质发展中心（学信网/研招网）}
\ith{学科差异}
\ist{认知与规划}
\begin{qa}{M0050}{保研要考政治英语吗？有没有笔试}
\textbf{不考统考的政治和外语，那正是免试免掉的部分。但要参加接收单位的复试，普遍有英语口语和专业能力考核；校内遴选环节多数学校明确规定不另设笔试。}

分三层说明：

\begin{enumerate}
  \item \textbf{统考初试科目免考}：教学〔2006〕14号第三条明确，免试是指应届本科毕业生不必经过全国硕士研究生入学统一考试的初试，直接进入复试。思想政治理论、外国语等初试科目因此不需要参加。
  \item \textbf{校内遴选通常不组织专门考试}：成都大学商学院2027届细则写明确保学业综合成绩作为遴选推荐综合评价体系的基础指标，除规定的专项计划外，不组织专门的遴选推荐考试（包括笔试、面试等）；广东外语外贸大学2026年办法同样规定不得专门组织遴选推免生的考试（包括笔试、面试等）。教学〔2006〕14号第十五条第（三）项对推荐环节的表述也只是学校按规定进行综合测评、择优确定初选名单，未要求组织考试。
  \item \textbf{接收单位的复试要考}：中南大学2027年选拔推免生直接攻博的通知写明复试内容包括外语水平、专业能力和综合素质及能力考核，并进行外国语的听力和口语测试；浙江大学2027年接收外校推免生通知规定复试成绩按百分制计算，不及格者（低于60分）不予录取，且须综合面试、全程录音录像。
\end{enumerate}

所以备考重点不在政治，而在英语口语、专业基础和面试表达。
\nref{S0454}
\nref{S0079}
\nref{S0134}
\nref{S0625}
\nref{S0032}
\end{qa}

\iin{教育部学生服务与素质发展中心（学信网/研招网）}
\ith{考核形式}
\ist{认知与规划}
\section{保后安排}
\begin{qa}{M0034}{保研一定要提前联系导师吗？不联系会怎样}
\textbf{取决于你报的是硕士还是直博，以及目标院校的安排。硕士阶段很多学校在录取后统一组织师生双选，不要求复试前联系；直博则常要求报名前先联系导师。}

三种不同的规定：

\begin{enumerate}
  \item \textbf{复试前无须联系、录取后统一双选}：北京航空航天大学仪器科学与光电工程学院接收推荐免试攻读2027年研究生工作方案写明，申请推荐免试攻读硕士研究生的考生参加复试前无须联系导师，学院将在复试结束后统一组织师生双选工作。
  \item \textbf{录取后互选、可提前进组}：哈尔滨工业大学计算学部关于接收2027年推免生的补充报名通知写明，推免生在录取后可与优秀导师进行互选，提前接受导师指导并进入课题研究工作，互选工作流程在录取后进行通知。
  \item \textbf{直博生须提前联系}：吉林大学生命科学学院2027年通知写明，招收直博生是指选拔具有推免生资格的优秀应届本科毕业生直接取得博士生入学资格的招生方式，直博生基础学制5年，并明确申请直博生的学生需在报名前与学院研办老师取得联系，以确定招生条件、范围及相关要求。
\end{enumerate}

所以不联系会怎样没有统一答案：硕士阶段可能完全没影响，直博阶段可能连报名资格都确认不了。先看目标院系通知里有没有硬性要求。
\nref{S0421}
\nref{S0552}
\nref{S0012}
\end{qa}

\iin{北京航空航天大学}
\ith{联系导师}
\ist{认知与规划}
\begin{qa}{M0044}{保研后大四是不是就轻松了？还要修课吗}
\textbf{不轻松。推免录取是有条件录取，大四这一年仍有几件必须完成的事，任何一件出问题都可能被取消录取资格。}

录取资格会被取消的常见情形（各校表述高度一致）：

\begin{itemize}
  \item 浙江大学2027年接收外校推免生通知列出五项：思想政治素质和品德考核不合格或政审不合格；从推免生报名至入学报到之日受到任何处分；本科最后一学年计入毕业要求学分的相关课程成绩出现不及格记录；本科毕业设计（论文）成绩未达到良好及以上（或80分及以上）；规定的研究生新生入学日期前无法提供毕业证书和学士学位证书。
  \item 中山大学2027年办法列出：未完成本科培养方案规定的所有课程及实践环节（含毕业论文或实习）的学分要求；本科最后一学年课程成绩出现不及格记录；有违法违纪行为或受到处分；政治表现不合格、体检不合格；不能按时取得毕业资格或取得毕业资格但未获得学士学位。
  \item 全国层面的兜底条款是教学〔2006〕14号第二十条：入学前未取得学士学位或本科毕业证书，或受到处分的，取消录取资格。
\end{itemize}

另有行为约束：北京大学2027届办法规定已落实接收单位的推免生不得再申请出国（境）留学、参加就业派遣等。
\nref{S0032}
\nref{S0167}
\nref{S0454}
\nref{S0553}
\end{qa}

\iin{浙江大学}
\ith{保后安排}
\ist{认知与规划}
\begin{qa}{M0053}{我现在大四了，九推还有机会吗？是不是都招满了}
\textbf{还有机会，但窗口很窄。九推指的就是推免服务系统开放后的正式填报阶段，2027届填报期是2026年9月21日9:00到10月18日，待录取确认从9月24日14:00开始。}

关键事实：

\begin{itemize}
  \item \textbf{系统时间}：2026年9月18日起注册并查看本人推免资格；9月21日9:00开始填报志愿；10月18日截止填报志愿；10月20日前未落实接收单位的推免生不再保留推免生资格（四川省教育考试院2027年公告）。
  \item \textbf{志愿可以多轮投}：可同时填报3个平行志愿，每个志愿提交后48小时内不允许修改；若招生单位拒绝或48小时后仍无复试通知、复试未通过，该志愿自动解锁，可继续填报其他志愿。
  \item \textbf{热门院校确实可能已满}：这不是制度限制，而是招生单位在夏令营和预推免阶段已经锁定生源。是否还有缺额，只能查目标院校研究生院公布的推免接收专业目录和当年的补录通知。
\end{itemize}

所以大四才决定走九推，务实策略是放弃只盯头部院校，把重点放在仍有缺额的院系，同时准备好万一落空就转统考（10月20日之后仍可报考当年统考）。
\nref{S0389}
\nref{S0182}
\nref{S0352}
\end{qa}

\iin{www.sceea.cn}
\ith{九推}
\ist{认知与规划}
\begin{qa}{M0058}{保研要不要提前学研究生课程？有意义吗}
\textbf{有意义，但不是必须。全国公开文件中没有推免生须提前修研究生课程的规定；能查到的相关说法，是部分院系允许录取后提前与导师互选、进入课题研究。}

可查到的做法是提前进组而不是提前修课。哈尔滨工业大学计算学部接收2027年推免生的补充报名通知写明，推免生在录取后可与优秀导师进行互选，提前接受导师指导并进入课题研究工作，互选工作流程在录取后进行通知。该学部同时说明，学术学位硕士研究生纳入硕博贯通培养体系，在硕士在读阶段可提前修读博士阶段相关课程；如申请硕博连读，最早可于研一下学期起开始申请。

需要注意，推免录取是有条件录取：浙江大学2027年通知明确，本科最后一学年计入毕业要求学分的相关课程成绩出现不及格记录、或毕业设计（论文）成绩未达到良好及以上，会被取消录取资格；中山大学2027年办法也要求完成培养方案规定的所有课程及实践环节的学分要求。

所以优先级应该是先把本科最后一年的课程和毕业设计做扎实，再考虑提前进组或自学研究生基础课。
\nref{S0552}
\nref{S0032}
\nref{S0167}
\end{qa}

\iin{哈尔滨工业大学}
\ith{保后安排}
\ist{认知与规划}
\section{排名口径}
\begin{qa}{M0016}{我们学校保研要笔试吗？还是只看排名}
\textbf{以本校细则为准。全国文件只要求推荐环节进行综合测评、择优确定初选名单，没有规定必须笔试；不少学校明确写不组织专门的遴选推荐考试。}

全国层面：教学〔2006〕14号第十五条第（三）项规定推荐程序为学校按照规定进行综合测评，择优确定初选名单，通篇没有要求组织笔试或面试。

院校层面出现两种明确做法：

\begin{itemize}
  \item \textbf{不另设考试}：成都大学商学院2027届细则写明，把学业综合成绩作为遴选推荐综合评价体系的基础指标，除规定的专项计划外，不组织专门的遴选推荐考试（包括笔试、面试等）；广东外语外贸大学2026年办法同样规定不得专门组织遴选推免生的考试。
  \item \textbf{可能有考核环节}：部分院校在综合评价中设专家审核小组鉴定加分材料。安徽师范大学2027年工作方案要求学院成立不少于7人的专家小组，对本院申请推免学生提交的科研论文、竞赛获奖、发明专利授权、入伍建功、国际交流及志愿服务等专长情况进行审核鉴定；云南师范大学2027年通知也要求组建推免工作监督小组全程监督并受理异议。
\end{itemize}

所以只看排名和要考试两种说法都可能成立。查法很简单：找本院当年推免细则，看遴选方式或综合成绩计算那一节。
\nref{S0454}
\nref{S0079}
\nref{S0134}
\nref{S0646}
\nref{S0568}
\end{qa}

\iin{教育部学生服务与素质发展中心（学信网/研招网）}
\ith{校内选拔}
\ist{认知与规划}
\begin{qa}{M0025}{退役大学生士兵专项计划是保研还是考研}
\textbf{要分成两类。退役大学生士兵专项硕士研究生招生计划是统考的专项招生计划，不是推免；而立功表彰退役军人免试生才是免初试通道，走推免服务系统报名。}

\textbf{一、退役大学生士兵计划（统考）}
教育部《2026年全国硕士研究生招生工作管理规定》第二十四条规定，退役大学生士兵计划：高校学生应征入伍退出现役且符合全国硕士研究生招生考试报考条件，可以申请报考。教育部办公厅《关于做好2026年退役大学生士兵专项硕士研究生招生计划招生工作的通知》（教学厅〔2025〕8号）明确该计划按“自愿报名、统一招考、自主划线、择优录取”的原则招生，2026年由北京大学、清华大学等532所普通高等学校承担；考生在报名时应按要求填报本人入伍前的入学信息以及入伍、退役等相关信息。此外，按教学〔2024〕4号《2025年全国硕士研究生招生工作管理规定》第二十四条，考生报名时须按要求填报信息，有关部门按职责分工审核；复试前须向招生单位提供《入伍批准书》和《退出现役证》供复核（最新年度表述以报考当年规定为准）。

\textbf{二、立功表彰退役军人免试生（免初试）}
同一管理规定第二十五条：对服役期间获得三等战功、二等功以上奖励或者二级以上表彰，符合全国硕士研究生招生考试报考条件的退役人员，可申请免初试攻读硕士研究生；符合免初试资格的考生应在教育部规定的全国统考报名时间内登录推免服务系统报名，逾期不得补报。四川省教育考试院2027年公告进一步明确，申请立功表彰免初试的考生不得再报名参加全国硕士研究生招生考试。

另外，推免名额不得为退役大学生士兵单独切块：云南师范大学2027年通知明确不得为仅符合单一或部分遴选条件的学生（如高水平运动员、退役大学生士兵、竞赛获奖者等）单列计划，但服役经历可作为综合评价的加分因素。
\nref{S0051}
\nref{S0052}
\nref{S0049}
\nref{S0389}
\nref{S0568}
\end{qa}

\iin{中华人民共和国教育部}
\ith{退役政策}
\ist{认知与规划}
\begin{qa}{M0033}{保研是看专业排名还是年级排名？小专业吃亏吗}
\textbf{多数学校按专业（或专业方向）排名，但全国文件没有统一规定，具体到哪一级、怎么算，全部由本校细则决定。}

全国层面的表述只到原则：教学〔2006〕14号第十五条第（三）项要求学校按照规定进行综合测评、择优确定初选名单，不涉及排名范围。院校则各自规定：

\begin{itemize}
  \item \textbf{排名口径由院系自定}：北京大学2027届办法写明学生成绩数据以教务系统统计为准，各院系的综合排名标准由院系自行制定和解释。
  \item \textbf{门槛写在本专业层面}：广东外语外贸大学2026年办法以获得本专业本科人才培养方案应修总学分的80\%、平均学分绩点不低于3.0（含3.0）、无课程不及格记录作为学业基本要求。
  \item \textbf{名额按专业切分}：成都大学商学院2027届细则规定，根据学校下达的年度推免生名额，结合学院各本科专业应届毕业班人数、本科一流专业建设及所对应的硕士学位点建设、重点学科建设需要，确定各专业的推免生名额及候补人选名额。
\end{itemize}

小专业吃亏是否成立，取决于名额是按专业切还是按学院统一排。名额切到专业，小专业的名额可能很少；学院统一排序，小专业学生反而可能因为同专业竞争少而占优。看细则里名额分配和综合成绩排序两节就能判断。
\nref{S0454}
\nref{S0553}
\nref{S0134}
\nref{S0079}
\end{qa}

\iin{教育部学生服务与素质发展中心（学信网/研招网）}
\ith{排名口径}
\ist{认知与规划}
\begin{qa}{M0054}{保研看成绩还是综合测评？哪个占比高}
\textbf{全国政策的导向是学业综合成绩作基础指标、其他表现作加分，加分项必须与学业成绩一起排序，不能单列名额；具体权重各校自定。}

政策原文的落点是教学〔2006〕14号第六条：在对考生平时学习和科研能力综合测评的基础上，突出对创新精神、创新能力和专业能力倾向的考查。院校把它落到具体权重上：

\begin{itemize}
  \item \textbf{明确基础指标地位}：广东外语外贸大学2026年办法规定将本科阶段学业综合成绩作为推免工作最基础的遴选指标；北京大学2027届办法规定把学生本科应届毕业学年前各学年的学业综合成绩作为遴选推荐综合评价体系的基础指标，可将参军入伍服兵役、院系认可的志愿服务、到国际组织实习以及高水平科研创新成果、竞赛奖项等纳入综合评价体系，但不得单列名额、单独排序。
  \item \textbf{权重可查的实例}：大连理工大学商学院的推免成绩排序方法为推免成绩＝学习成绩×0.85＋思想素养＋综合素质和能力测评×0.1，其中学习成绩满分100分。
  \item \textbf{加分项要提前公布}：云南师范大学2027年通知要求各单位组织专家论证并提出本学科高水平科研创新成果、竞赛奖项等加分指标的认定范围、标准，并提前公布。
\end{itemize}

所以成绩还是综合测评这个问法的答案是：两者都在，但成绩通常是权重最大的那一块，且加分项受不得单列名额的限制。
\nref{S0454}
\nref{S0134}
\nref{S0553}
\nref{S0549}
\nref{S0568}
\end{qa}

\iin{教育部学生服务与素质发展中心（学信网/研招网）}
\ith{排名口径}
\ist{认知与规划}
